\section{Benchmark and evaluation families}\label{sec:evaluation}


\regime{Control}{what the benchmark literature measures about state, and what it leaves untested.}
If always-on agents are persistent-state systems rather than memory-augmented chat models, evaluation has to ask whether a benchmark measures governed state over a lifecycle, instead of only whether a system can store and retrieve memory. A benchmark is an operational definition: it fixes what counts as success, and therefore what counts as a system worth building. We read the evaluation literature against six recurring benchmark stressors: persistent multi-user memory, stale-state adjudication, social and belief dynamics, real action with side effects, privacy and consent boundaries, and recovery after bad memory. These stressors are not a separate theory of the paper; they are a compact checklist derived from the axes and lifecycle. Persistent multi-user memory stresses scope and provenance; stale-state adjudication stresses mutability and authority; social and belief dynamics stress shared state and provenance; real side effects stress actionability and recoverability; privacy and consent stress scope and authority; and recovery stresses rollback traceability and deletion propagation. The argument is structural rather than accusatory. Many benchmarks are excellent within scope, but their scopes divide the always-on evaluation space so cleanly that no family is forced to govern state, act on it, and recover from it in one protocol. Recoverability under authority and provenance constraints then falls between families.

Five benchmark families define the current measurement surface. Memory question-answering over histories made memory failures observable but evaluates state after the fact. Long-context and RAG benchmarks are boundary cases: they constrain what supplied context a model can use, but they are not memory benchmarks and were never meant to be. Memory-guided action benchmarks move closest to deployment by making earlier experience change later choices, yet they reset identity, consent, and governance at every task boundary. State and belief mutation benchmarks attack obsolete and conflicting state most directly and are where authority first becomes visible, but they rarely couple mutation to real action or deletion. Personalization and proactive benchmarks test long-horizon user state and, more recently, spontaneous recall, but treat write policy and consent as out of scope. Cutting across all five is a methodological literature on judge reliability and statistical rigor that determines whether any of these scores can be trusted at all. Methodology comes first as a lens; the family-by-family reading then closes with the coverage finding that motivates a new protocol.

\subsection{Evaluation methodology and what a score can mean}

Before comparing what each family measures, we ask what an agent benchmark score can support, because the always-on setting strains standard evaluation methodology in specific ways. A recurring critique is that agent evaluation conflates capability with measurement artifact. \citet{kapoor2024aiagentsthatmatter} document that agent benchmarks frequently omit cost-versus-accuracy reporting, lack adequate holdout sets so that systems overfit to the benchmark, and rarely report failure rates, all of which inflate apparent progress. \citet{miller2024errorbarsevals} supplies the missing statistical machinery, framing benchmark scores as estimates over a super-population with confidence intervals and principled two-model difference tests, and \citet{madaan2024quantifyingvariance} measures seed variance, training-step non-monotonicity, and continuous-metric instability empirically, showing that single-number leaderboard comparisons are often within noise. Contamination is a parallel threat: \citet{white2024livebench} responds with monthly-refreshed questions drawn from recent sources and objective ground truth so that memorized answers cannot masquerade as reasoning. These concerns sharpen in the always-on regime, where a system is scored on a trajectory of dependent decisions rather than independent items: errors compound across sessions, variance accumulates, and a single contaminated fact can be silently reused for the remainder of an episode.

A second methodological axis is the rise of LLM-as-judge and its known biases, which matters acutely once we score state transitions rather than final answers. \citet{zheng2023judgingllmasajudge} established the judge paradigm via MT-Bench and Chatbot Arena and reported strong agreement with human preference, but also documented position, verbosity, and self-enhancement biases. Subsequent work quantified and partially mitigated these effects: \citet{shi2024judgingthejudges} measures position bias in pairwise and list-wise judging with repetition-stability and position-consistency diagnostics; \citet{tripathi2025pairwiseorpointwise} shows that the protocol choice itself induces systematic bias, with roughly a third of preferences flipping under distractors depending on whether scoring is pointwise or pairwise; and \citet{lee2024checkeval} replaces unreliable Likert scoring with decomposed binary checklist questions to raise reliability. The judge-reliability surveys of \citet{gu2024surveyllmasajudge} and \citet{li2024llmsasjudgessurvey} synthesize consistency-improvement, bias-mitigation, and meta-evaluation methodology, while \citet{chiang2024chatbotarena} provides the crowdsourced pairwise plus Bradley-Terry ranking infrastructure that anchors much of the field's reliability evidence. \citet{you2026agentasjudge} extends the paradigm from judging a final answer to judging an agentic process step by step. Once the judge is itself an agentic process, its reliability becomes a moving target that interacts with the very governance policies under test: \citet{iacob2026redqueen} co-evolves the evaluation framework with the agent's governance objectives, detecting bias in agent-as-judge systems and showing that self-improvement guarantees must be re-established as authority and scope policies shift, which is exactly the regime an always-on benchmark scores. The methodological lesson is double-edged. Governance properties such as whether a deletion propagated or whether an action was licensed by a current authority epoch are exactly the multi-step, process-level judgments where a naive scalar judge is least reliable. Yet several governance checks are crisp predicates (was the tombstoned value retrievable, yes or no) that admit a deterministic oracle and need no judge at all. This split, between properties that demand fragile holistic judgment and those that reduce to deterministic invariant checks, recurs across the families and foreshadows why the protocol we ultimately motivate keeps its governance checks deterministic.

\subsection{Memory QA over histories}

The family that founded memory evaluation asks questions over long interaction histories. LoCoMo \citep{maharana2024evaluatinglongtermconversationalmemory} is the canonical instance: very long multi-session dialogues with probes that require recovering semantically distant facts, reasoning across sessions, handling temporal updates, abstaining when the answer is absent, and selectively forgetting. Its contribution was to make memory failure a measurable quantity rather than an anecdote, and it remains a standard reference point. The lineage is older than LoCoMo, and reading it clarifies what the family does and does not test. Multi-session open-domain dialogue datasets first established recall and summarization across sessions, and later work typed personal memory into semantic facts, episodic events, and social relations, on which LoCoMo-style benchmarks build. The shared assumption is that memory quality can be read off the correctness of answers to retrospective questions.

That assumption is precisely where the family draws methodological fire from within. \citet{flynt2026precisionmembench} argues that LoCoMo-style benchmarks conflate answer-correctness with retrieval-correctness: a system can answer correctly with the wrong evidence in context, or fail despite retrieving the right evidence, and pooled accuracy hides both. PrecisionMemBench responds by scoring retrieval precision in isolation from the generator, so that memory quality is no longer entangled with the language model's ability to paper over bad retrieval. \citet{long2026memtrace} pushes the same disaggregation further, re-basing long-term-memory evaluation on the per-fact knowledge point probed across age, question-type, and evidence conditions, and finds that the dominant bottleneck is not retrieval at all but evidence use: agents often have the right fact in context and still answer wrong. \citet{deng2026entitycollision} attacks a subtler confound, showing that apparent retrieval lift in agent memory can be lexical leakage rather than semantic competence; their stratified protocol pins a BM25 floor via entity-token collision so that improvements are attributable to the embedder. Together these works are a corrective lineage: each shows that a higher LoCoMo-style number can be produced without the underlying capability the number is supposed to certify.

This family therefore measures recall, temporal reasoning over updates, abstention, and, with the newer protocols, retrieval precision and evidence use cleanly separated from generation. It does not measure the write and governance half of the lifecycle. These benchmarks ask questions after the fact; they rarely test whether a remembered fact should have been written, whether the write was authorized, whether retrieving it would improve a later action, or what happens when the user later asks for it to be deleted. There is no notion of a permission epoch under which a fact was valid, no provenance tier distinguishing a user-stated preference from a fact scraped from an untrusted tool output, and no action whose side effect could be wrong because the recalled state was stale. The gap is intrinsic to the format: a question-answer pair is a recall unit, and a recall unit has no field for an authority tag, a deletion record, or a rollback trace. Even the most rigorous member of this family, having stripped out the retrieval-versus-generation confound, still certifies only that the right text can be recovered, not that recovering and acting on it is governed. Memory QA isolates retrieval fidelity and leaves downstream governance outside the test.

\subsection{Boundary cases: long-context and RAG benchmarks}

Long-context and retrieval-augmented-generation benchmarks are frequently cited as evidence about agent memory, but they are boundary cases rather than members of the family. They evaluate whether a model can use information already supplied to it, not whether a system should commit information to durable agent-owned state and govern it thereafter. LongBench \citep{bai2024longbenchbilingualmultitaskbenchmark} probes long-context understanding across diverse bilingual tasks and establishes an empirical baseline for using a large window; \citet{liu2023lostmiddlelanguagemodels} shows the sharp boundary effect that models underuse information placed in the middle of a long context despite attending to its edges. On the RAG side, \citet{chen2023benchmarkinglargelanguagemodels} decomposes retrieval quality into noise robustness, negative rejection, information integration, and counterfactual handling, giving the field a vocabulary for faithfulness and robustness to noisy evidence. \citet{gutierrez2024hipporag} bridges toward structured memory by organizing retrieval over a knowledge graph with a PageRank-style spreading-activation step, improving multi-hop retrieval, but its graph is built offline and queried statically.

These benchmarks measure context use, retrieval faithfulness, and noise robustness extremely well, and the always-on field should treat their results as binding constraints: if a model cannot reliably use the middle of its context window, no memory system that surfaces evidence into that window is safe. But the limitation is structural and definitional. A long-context benchmark holds the information fixed and varies the model; an always-on system must decide what becomes durable state, when that state becomes authoritative, and how it is later revised, scoped, revoked, or removed. RAG benchmarks score retrieval against a fixed corpus; they do not score a durable write, an update under a temporal-validity window, a deletion that must propagate to derived tiers, or an authority constraint on who may read which evidence. The boundary effect documented by lost-in-the-middle is a property of inference over a static window; the always-on failure mode is that the window is assembled from accumulated state whose provenance, freshness, and authorization are themselves uncertain. This family therefore brackets the substrate problem rather than the governance problem. Its measures of durable write, update, deletion, and authority are absent, by design and correctly so, because that was never its question. Treating long-context or RAG scores as memory-governance scores is a category error, and recognizing the boundary is itself part of the argument: the field has mature instruments for the supply-and-use of context and almost none for the governance of state.

\subsection{Memory-guided action benchmarks}

The family that moves closest to always-on deployment makes earlier experience change later choices, so that memory is evaluated through its effect on action rather than through retrospective questions. MemoryArena \citep{he2026memoryarena} is a clear recent example: tasks are interdependent across sessions, so an action or feedback in one session should alter behavior in a later one, and the benchmark reports recall, induction, and contextual grounding. LifelongAgentBench \citep{zheng2025lifelongagentbenchevaluatingllmagents} evaluates agents as lifelong learners over sequential task streams, measuring forward and backward transfer and catastrophic forgetting in serial tool environments. Surrounding these are the realistic task substrates that show how hard action remains even before persistent personalization is added: enterprise knowledge-work tasks such as WorkArena \citep{boisvert2024workarena}, open-ended computer-use environments such as OSWorld \citep{xie2024osworldbenchmarkingmultimodalagents}, and stateful application benchmarks such as AppWorld \citep{trivedi-etal-2024-appworld}, where an agent must track world-state changes across interactions with simulated apps and people. This family's contribution is the move from final-answer scoring to final-state checking: success is verified against the state of the world after the agent acts, which is the right test for a system whose state has actionability.

Comparing these benchmarks exposes the shared limitation. MemoryArena and LifelongAgentBench reward experience reuse but treat memory as an unaudited performance booster: there is no scoring of whether a reused experience was authorized to influence the new task, no provenance tier on the stored trajectory, and no test of whether a poisoned or stale experience should have been blocked. The realistic substrates are, almost by construction, episodic: WorkArena, OSWorld, and AppWorld reset the environment at each task, which is exactly the boundary an always-on protocol must cross. They verify that an action produced the right final state, but reset identity, consent, accumulated permission, and long-run governance between tasks, so the questions that define always-on behavior never arise. Even AppWorld, which executes real side effects against stateful apps, scores whether the side effect was correct, not whether it was recoverable: there is no perturbation that retries an irreversible call to test idempotency, no revocation of the authority that licensed an earlier action followed by a check that the action's state was rolled back, and no deletion request whose propagation to derived state is verified. The matrix view makes the pattern concrete: each suite tracks state in action but does not isolate whether failures stem from a bad write, weak organization, or weak retrieval at action time, and none couples action to governed writes, revocation, or rollback. The gap here is not realism and not action, both of which this family has in abundance. The gap is that the same protocol does not jointly stress persistent identity, governed writes, provenance, deletion, rollback, stale-state adjudication, privacy scope, and recovery. Memory-guided action benchmarks prove that acting on memory is hard and measurable; they leave untested whether acting on memory is governed.

\subsection{State and belief mutation benchmarks}

The fourth family attacks the problem the others defer: what happens when stored state becomes obsolete, conflicting, or contested, and which stored state should govern the present. This is where authority first becomes an explicitly measured property rather than an implicit assumption. STALE \citep{chao2026stalellmagentsknow} asks whether agents detect that a stored memory has gone stale and can resolve an implicit conflict between an obsolete value and a current one, probing the validation stage of the lifecycle. \citet{han2026memorydepth} approaches the same staleness-and-expiry question from the write side rather than the read side, evaluating whether a long-running agent correctly preserves still-valid facts and expires obsolete ones across context teardown by gating consolidation on surprise and valence signals, which makes \emph{what state survives} a scored property rather than an implementation detail. In the social direction, GroupMemBench \citep{yang2026groupmembenchbenchmarkingllmagent} is the first benchmark for agent memory in multi-party conversation, testing speaker attribution and turn-by-turn state tracking so that who said what, and therefore whose authority a remembered fact carries, becomes scoreable; and the interactive social-intelligence setting of SOTOPIA \citep{zhou2024sotopiainteractiveevaluationsocial} reveals how persistent memory underwrites believable multi-turn social behavior. A cluster of newer benchmarks push on individual always-on stressors, each isolating one axis. \citet{jung2026meme} is a direct probe of deletion propagation: its multi-entity evolving-memory tasks, Cascade, Absence, and Deletion, test whether an update or removal propagates to dependent entries, and the reported near-total collapse (Cascade at 3 percent, Absence at 1 percent) is evidence that tested memory systems struggle to maintain the deletion-propagation invariant. \citet{kwon2026reclaim} isolates a provenance failure with a striking consequence: a memory that keeps a stale conclusion but drops its source becomes confidently uncorrectable, strictly worse than empty memory, and the judge-free reclaim evaluation shows a source-first write policy repairs it. \citet{zhang2026kpr} stress-tests the converse condition, whether an agent preserves provenance and re-verifies authority before it accepts external knowledge, using controlled ablation of descriptions and diffs together with poisoned-patch injection, so that provenance retention and authority re-checking on ingestion become directly measurable rather than assumed. \citet{wang2026docarena} attacks scope enforcement at benchmark-construction time, automatically translating document access policies into adversarially structured question-answer environments that test whether an agent respects tool-use scope and can verify the provenance chain of the evidence it returns, turning an access-control specification into a graded test of scoped retrieval. \citet{bhargava2026memaudit} audits what memory writers actually preserve under a fixed storage budget, isolating selection quality and validity-state preservation from raw representation, and \citet{shao2026scaleconditioned} holds task evidence fixed while adding irrelevant sessions to decompose where stored evidence stops being usable under scale. \citet{liu2026streammembench} brings this into a streaming, egocentric setting and finds that systems often fail to turn stored evidence and feedback into reliable future behavior, while \citet{zheng2026seagym} scores self-evolving harnesses by decomposing an update into reusable gain, overfit, regression, and cost, revealing that an intermediate self-improvement snapshot can collapse later. Underlying the whole family is a formal claim: \citet{orogat2026gem} reframes long-term memory correctness as a property of the state trajectory rather than of individual records, defining Governed Evolving Memory with state-level operators and six correctness conditions covering write, read, delete, consistency, and rollback traceability, and proving that record-level stores cannot satisfy them.

This family is the most important for the survey's thesis, because it is the only one in which authority, staleness, provenance, and deletion are first-class measured quantities rather than incidental. It is also where the field's governance deficit shows up as hard numbers: the MEME collapse and the RECLAIM uncorrectability result are not warnings about a hypothetical future, they are measurements that today's stores fail invariants the thesis names. Yet the coupling gap remains. These benchmarks test mutation largely in isolation from action: STALE asks whether the agent knows a value is stale, not whether it refrained from issuing a tool call licensed by that stale value; MEME tests whether a deletion propagated through the memory graph, not whether a side effect derived from the deleted value was rolled back; GEM formalizes rollback traceability as a correctness condition but is a framework awaiting an empirical harness that audits real systems against it. The social members add belief and group structure but, as the matrix notes, do not score authority and scope governance during multi-party updates or deletion propagation when a speaker leaves the group. Even when a benchmark records an authority or scope violation, the violation must be interpreted: \citet{singh2026modelforensics} supplies a forensic baseline that combines chain-of-thought inspection with causal intervention to distinguish a genuine governance breach from a harmless reasoning shortcut, the diagnostic any always-on protocol needs before it can attribute a failed invariant check to real misalignment rather than a benign path. So the family that most directly tests state mutation still rarely couples mutation to real tool action, to a privacy or consent boundary, or to recovery after a bad write. It proves that stale and contested state is mishandled; it does not prove whether a system can act on contested state and then recover.

\subsection{Personalization and proactive benchmarks}

A fifth family evaluates long-horizon user state and, increasingly, whether an agent surfaces stored state on its own initiative. The personalization line operationalizes what it means for state to be long-term and user-owned. \citet{xu2026personalizedllmpoweredagentsfoundations} surveys evaluation criteria for personalized LLM agents, defining multi-session coherence and preference retrieval as the targets; \citet{jiang2025personamem} contributes PersonaMem, which tracks an evolving user persona across interactions and reports persona-recall metrics; and \citet{xie2026dynamicmem} sharpens the temporal dimension by separating user attributes, habits, and preferences that mutate on different timescales, exposing that user state evolves at multiple speeds and that a single freshness policy cannot fit all of them. On the action-coupled side, \citet{cai2025personalizedwebagents} demonstrates with PUMA and the PersonalWAB benchmark that retrieving user history to align cross-session web actions measurably helps, closing part of the loop from personalization to action that the memory-QA family leaves open. The proactive line is newer and addresses a stressor no other family touches: whether an agent recalls and acts on a latent stored constraint without being prompted to. \citet{zhang2026triggerbench} measures this prospective memory, finding that spontaneous recall is far harder than retrospective query-driven recall and degrades as context length grows. This is a distinct competence: every other benchmark cues the relevant memory through the question or the task, whereas a deployed always-on agent must decide on its own when stored state becomes relevant and when surfacing or acting on it is appropriate.

The contrast within this family is instructive. Personalization benchmarks measure persona recall, multi-session coherence, and preference adaptation well, and PUMA shows personalization can be tied to action; TriggerBench adds the orthogonal and under-measured competence of utility-calibrated proactivity. But the governance half is again out of frame. As the matrix records, PersonaMem tests recall but not recoverability or auditability, with no validation that an agent can undo or correct a poisoned persona fact or detect a contradiction; PUMA is retrieval-based with no scope control over which user facts may influence which actions and no audit on action provenance; and the dynamic-attribute benchmark tests mutability but not provenance, leaving no record of which fact changed, when, or which user action caused the change. TriggerBench measures whether a constraint is recalled in time, not whether the agent distinguishes a constraint it is authorized to act on from one whose authority has lapsed, nor whether a proactive action is recoverable if the triggering memory was stale or poisoned. Personalization is where consent and privacy should be most salient, since the stored state is by definition sensitive and user-owned, yet write policy, consent, and the right to erasure are treated as out of scope. The clearest move to close this is to make the substrate itself governance-aware: \citet{selvam2026profilefoundry} builds a synthetic, deterministic profile substrate with generational provenance so that private-fact leakage, tool-use scope enforcement, and consistent state recovery across invocations can be scored as first-class outcomes, showing that the consent and recoverability dimensions are testable in personalization once the data is generated with provenance and scope baked in rather than retrofitted onto a recall set. The family extends always-on evaluation in two directions, long-horizon user state and spontaneous recall, while leaving the same governance gap open: it tests that the agent remembers the user and can act on that memory, not that it is licensed to, can be told to forget, and can recover when the personalized state proves wrong.

\subsection{Coverage against the six always-on benchmark stressors}

Reading the five families together yields the main evaluation finding of this survey: their coverage of the six always-on stressors is a clean partition rather than a single missing dataset. Table~\ref{tab:eval-families} summarizes all five families qualitatively, by representative examples, what each measures well, and what it leaves untested for always-on agents. For three of the families, where the closest-work set is dense enough to support it, we additionally coded individual benchmarks against the six stressors at the per-benchmark level; the long-context/RAG and personalization families are characterized at the family level only, because their member benchmarks were not designed around these stressors and a per-benchmark coding would mostly record absences.

The partition is visible at the level of individual benchmarks as well as at the family level. Table~\ref{tab:eval-coverage} codes representative benchmarks from the three families where the survey's benchmark-level coding is densest against the six stressors. Reading down the columns is more telling than reading across the rows: the recovery column (R6) is almost entirely empty, and no single benchmark scores real action (R4), a privacy boundary (R5), and recovery (R6) together, which is the combination an always-on agent most needs and the one the field most consistently leaves unmeasured.

\begin{table}[t]
\centering\small
\setlength{\tabcolsep}{4pt}\renewcommand{\arraystretch}{1.15}
\begin{tabularx}{\linewidth}{p{0.20\linewidth}Ycccccc}
\toprule
Benchmark & Family & R1 & R2 & R3 & R4 & R5 & R6 \\
\midrule
LoCoMo & Memory QA &  & \pmark &  &  &  &  \\
LongMemEval & Memory QA &  & \pmark &  &  &  &  \\
MemoryAgentBench & Memory QA &  & \pmark &  & \pmark &  &  \\
PerLTQA & Memory QA &  & \pmark & \pmark &  &  &  \\
MemoryArena & Mem-guided action &  & \pmark &  & \cmark &  &  \\
LifelongAgentBench & Mem-guided action &  & \pmark &  & \cmark &  &  \\
WorkArena & Mem-guided action &  &  &  & \cmark &  &  \\
WorkArena++ & Mem-guided action &  &  &  & \cmark &  &  \\
STALE & State/belief mut. &  & \cmark &  &  &  &  \\
Memora/FAMA & State/belief mut. &  & \cmark &  &  &  &  \\
BeliefShift & State/belief mut. &  & \cmark & \cmark &  &  &  \\
Lifelong-SOTOPIA & State/belief mut. & \pmark &  & \cmark &  &  &  \\
\bottomrule
\end{tabularx}
\caption{Per-benchmark coverage of the six always-on benchmark stressors, for representative benchmarks from the three families where benchmark-level coding is densest. R1 persistent multi-user memory, R2 stale-state adjudication, R3 social and belief dynamics, R4 real action with side effects, R5 privacy or consent boundary, R6 recovery after bad memory. $\bullet$ denotes a scored stressor, $\circ$ partial, blank not scored. Two patterns are visible by column: the R6 (recovery) column is essentially empty, and no row marks R4, R5, and R6 together. The partition is structural, not a single missing dataset.}
\label{tab:eval-coverage}
\end{table}


\begin{table}[t]
\centering
\small
\setlength{\tabcolsep}{4pt}
\renewcommand{\arraystretch}{1.2}
\begin{tabularx}{\linewidth}{p{0.135\linewidth}YYY}
\toprule
Family & Representative examples & Measures well & Missing for always-on \\
\midrule
Memory QA over histories & LoCoMo \citep{maharana2024evaluatinglongtermconversationalmemory}; PrecisionMemBench \citep{flynt2026precisionmembench}; MemTrace \citep{long2026memtrace} & Recall, temporal updates, abstention, retrieval precision vs.\ generation & Write policy, authority, consent, real action, recovery \\
Boundary: long-context and RAG & LongBench \citep{bai2024longbenchbilingualmultitaskbenchmark}; Lost-in-the-Middle \citep{liu2023lostmiddlelanguagemodels}; RGB \citep{chen2023benchmarkinglargelanguagemodels} & Context use, retrieval faithfulness, noise robustness & Durable write, update, deletion, authority (out of scope by design) \\
Memory-guided action & MemoryArena \citep{he2026memoryarena}; LifelongAgentBench \citep{zheng2025lifelongagentbenchevaluatingllmagents}; AppWorld \citep{trivedi-etal-2024-appworld}; OSWorld \citep{xie2024osworldbenchmarkingmultimodalagents} & Experience reuse, tool and computer use, final-state checking & Persistent identity, governed writes, revocation, rollback \\
State and belief mutation & STALE \citep{chao2026stalellmagentsknow}; MEME \citep{jung2026meme}; RECLAIM \citep{kwon2026reclaim}; GroupMemBench \citep{yang2026groupmembenchbenchmarkingllmagent} & Staleness, deletion propagation, provenance, belief and group dynamics & Coupling to real action, privacy boundary, recovery after bad write \\
Tool-security evaluation & AgentDojo \citep{debenedetti2024agentdojodynamicenvironmentevaluate}; AgentPoison \citep{chen2024agentpoisonredteamingllmagents} & Prompt-injection attacks and defenses over realistic tool/data boundaries & Persistent write, deletion propagation, rollback after injected state persists \\
Personal\-ization and proactive & PersonaMem \citep{jiang2025personamem}; PersonalWAB \citep{cai2025personalizedwebagents}; TriggerBench \citep{zhang2026triggerbench} & Persona recall, preference drift, spontaneous proactive recall & Consent and erasure, scope control, audit, recoverability \\
\bottomrule
\end{tabularx}
\caption{Benchmark and evaluation families relevant to always-on agents. Each family measures part of the requirement space well and leaves a disjoint part untested. No family is forced to govern state, act on it, and recover from it at once, so recoverability under authority and provenance constraints falls between families.}
\label{tab:eval-families}
\end{table}

The six stressors an always-on benchmark would need to exercise are persistent multi-user memory, stale-state adjudication, social and belief dynamics, real tool or GUI action with side effects, a privacy or consent boundary, and recovery after a bad memory. The coverage finding is that, to our knowledge, no widely adopted benchmark we examined covers more than five of the six, and, more pointedly, none combines real action, recovery, and a privacy or consent boundary in a single protocol. The reason is visible in the table: the stressors are distributed across families so that satisfying all six would require a benchmark to inherit the action realism of the memory-guided family, the mutation adjudication of the state-and-belief family, the consent salience of the personalization family, and a recovery mechanism that essentially no family scores. Recovery is the binding stressor. As the survey's corpus statistics show, only a small fraction of works expose any rollback mechanism, and none of the benchmark families reports recovery success or cost after corruption, so even where a benchmark perturbs state it cannot certify that a system restored a known-good state afterward. Action without recovery is the most dangerous combination to leave unmeasured, because an always-on agent that acts on memory and cannot recover from a bad write is precisely the failure the thesis warns about, and it is exactly the combination that the action family executes but does not test and the mutation family tests but does not execute.

The deficit is therefore structural rather than a gap that one more dataset would close. Each family's format encodes its blind spot. A memory-QA recall unit has no field for an authority tag or a deletion record; an episodic action benchmark resets the very persistence it would need to test governance across; a mutation benchmark scores the memory graph but not the side effect derived from it; a personalization benchmark treats the user's stored state as a recall target rather than as consent-bound, revocable, auditable state. None of these is a defect of the benchmark within its own purpose. The defect is that the union of their purposes still leaves the persistent-state lifecycle, observe, write, validate, organize, retrieve, act, update, forget, audit, rollback, evaluated only in pieces, with the governance-defining stages of forget, audit, and rollback measured least and never jointly with action. A score on any one family certifies a capability that a stateless or single-session system could in principle also possess; none certifies the property that actually distinguishes a persistent-state system, namely that state is governed across its lifecycle under authority, scope, provenance, recoverability, and actionability constraints.

The evaluation gap motivates a separate protocol. The missing instrument is not a harder recall set or a larger action environment, but one that records and scores state transitions rather than answer quality: an episode whose stream carries idempotency keys, causal links, permission epochs, provenance tiers, and retention constraints, and whose perturbations include restart, conflicting update, deletion request, adversarial write, shared-scope change, permission revocation, and delayed consequence. The score should come from deterministic invariant checks of the kind the methodology subsection argued these predicates admit, not from a bias-prone holistic judge. Such an instrument would treat a trajectory as correct only if a deleted value is truly unretrievable across derived tiers, an action is licensed by a current rather than a lapsed authority epoch, a conflict is surfaced with provenance, a retried call does not duplicate an irreversible side effect, and a rollback trace is sufficient to restore a known-good state. Because the five families partition these stressors rather than combining them, the next section develops the Always-On Evaluation Protocol: an event-stream and snapshot schema with worked episodes and deterministic checks, designed to wrap the realistic substrates this family already provides with the restart, conflict, revocation, leakage, and rollback perturbations that govern persistent state.
